% =============================================================================
% Kapitel 6: Verifikation und Validierung
% =============================================================================

\section{Verifikation und Validierung}
\label{sec:verification_validation}

% =============================================================================
\subsection{Zielsetzung und Evidenzstrategie}
\label{sec:vv_objective}
% =============================================================================

Die in Kapitel~4 entwickelte Informationsverarbeitungsarchitektur und ihre in
Kapitel~5 beschriebene prototypische Realisierung werden im Folgenden anhand
mehrerer komplementärer Evidenzformen untersucht. Ziel ist nicht, aus einem
einzelnen Testtyp auf die Gesamtqualität des Systems zu schließen, sondern
unterschiedliche technische und fachliche Prüfgegenstände getrennt zu
betrachten und anschließend zu einer belastbaren Evidenzbasis für den
Proof of Concept zusammenzuführen.

Die Prüfgegenstände werden dabei aus der in Kapitel~4 beschriebenen
Architecture Description und den dort adressierten Stakeholder Concerns
abgeleitet. Die begriffliche Einordnung von Architecture Description,
Stakeholder und Concern orientiert sich an ISO/IEC/IEEE~42010:2022
\cite{ISO42010_2022}. Damit wird die Verifikation und Validierung nicht
losgelöst vom Architekturmodell definiert, sondern auf diejenigen
Verantwortlichkeiten, Informationsgrenzen und Interaktionen bezogen, die
zuvor als Teil der Architektur beschrieben wurden.

Die Verifikations- und Validierungsstrategie folgt damit keinem einzelnen
abschließenden Akzeptanztest. Sie kombiniert automatisierte Contract-,
Integrations- und Regressionstests mit kontrollierten End-to-End-Durchläufen,
einer eigenständigen Multi-Source-Verifikation, externer
SysML-v2-Validierung sowie expliziten menschlichen Review- und
Freigabeschritten.

Diese Kombination entspricht dem in Kapitel~3 beschriebenen iterativen
Untersuchungsansatz. Die prototypische Realisierung dient dabei nicht nur als
technischer Demonstrator, sondern als Untersuchungsartefakt, an dem
Architekturannahmen, Informationsgrenzen und Authority Contracts schrittweise
überprüft werden können. Die Evaluation eines entwickelten Artefakts ist dabei
ein wesentlicher Bestandteil designorientierter Forschung
\cite{Wieringa2014}.

Die einzelnen Evidenzformen besitzen unterschiedliche Aussagegrenzen.
Automatisierte Tests können beispielsweise die Einhaltung technischer
Contracts oder die Reproduzierbarkeit deterministischer Verarbeitung
überprüfen. Sie können jedoch keine fachliche Engineering-Freigabe ersetzen.
Umgekehrt kann eine menschliche Bewertung der fachlichen Plausibilität eines
Modellergebnisses keine technische Integritäts- oder Regressionsprüfung
ersetzen. Kapitel~6 behandelt diese Evidenzformen deshalb getrennt, bevor
Kapitel~7 die daraus resultierenden Beobachtungen zusammenführt.

\subsubsection{Pruefgegenstaende und Claim Boundaries}
\label{sec:verification_subjects}

Die Verifikation und Validierung adressiert nicht ausschließlich das erzeugte
SysML-v2-Artefakt am Ende der Verarbeitung. Prüfgegenstand ist vielmehr die
gesamte in Kapitel~5 beschriebene Kette der Informationsverarbeitung mit ihren
jeweiligen Übergängen und Authority Boundaries.

Auf technischer Ebene umfasst dies zunächst die verwendeten
Datenstrukturen, Identitäten, Manifeste, Fingerprints und
Verarbeitungscontracts. Von besonderer Bedeutung sind dabei diejenigen
Regeln, die ungültige, unvollständige oder veraltete Artefaktzustände
fail-closed von einer weiteren Verarbeitung ausschließen.

Darauf aufbauend werden die Übergänge zwischen den wesentlichen
Verarbeitungsstufen betrachtet. Dies betrifft insbesondere die Bindung von
source-grounded Evidence an Engineering Subjects, die Promotion zu Approved
Engineering Information, die Zulassung durch Project Fit, die Ableitung und
menschliche Platzierung von Model Candidates sowie die weitere Authority Chain
über Model Assembly Review, Base Internal Engineering Model,
Target-Model Formulation, Model Quality Review und Successor Internal
Engineering Model.

Ein weiterer eigenständiger Prüfgegenstand ist die Multi-Source-Verarbeitung.
Mehrere erfolgreich durchlaufene Single-Source-Pfade reichen hierfür nicht
aus. Die Multi-Source-Fähigkeit erfordert den Nachweis, dass mehrere Sources
innerhalb eines gemeinsamen Project-Kontexts verarbeitet werden können, ohne
Source Provenance, Artifact Identity oder die bestehenden Human-Authority-
Grenzen zu verlieren.

Für den erzeugten Zielzustand werden schließlich die deterministische
SysML-v2-Generierung, die interne Generated-Artifact-Prüfung, die externe
Validierung mit SYSIDE sowie Final Human Model Review und kontrollierte
Publication betrachtet.

Die daraus ableitbaren Claims bleiben bewusst auf den untersuchten
Proof-of-Concept-Scope begrenzt. Die V\&V kann Evidenz dafür liefern, dass die
konzipierte Architektur technisch realisierbar ist, zentrale Contracts und
Authority Boundaries eingehalten werden und der implementierte
Informationsfluss für die untersuchten Single-Source- und Multi-Source-Fälle
bis zu einem validierbaren SysML-v2-Artefakt durchlaufen werden kann.

Nicht daraus abgeleitet werden können eine universelle fachliche Korrektheit
erzeugter Engineering-Modelle, eine vollständige Abdeckung von SysML v2,
statistische Aussagen zur Usability, unbegrenzte Skalierbarkeit oder
Production Readiness. Diese Begrenzungen werden in Kapitel~8 im Kontext der
Forschungsfragen erneut aufgegriffen.

\subsubsection{Technische und fachliche Evidenz}
\label{sec:technical_vs_engineering_evidence}

Für die Bewertung des Proof of Concept werden technische und fachliche
Evidenz bewusst voneinander getrennt.

Technische Evidenz entsteht insbesondere durch automatisierte Unit-,
Contract-, Integrations- und Regressionstests sowie durch deterministische
Validierungen von Artefaktzuständen und deren Bindungen. Diese Prüfungen
adressieren beispielsweise gültige Identitäten, vollständige Authority
Contracts, konsistente Fingerprints, reproduzierbare Übergaben zwischen
Verarbeitungsschritten und die Einhaltung fail-closed definierter
Fehlerzustände.

Eine weitere technische Evidenzform bildet die externe Validierung des
generierten SysML-v2-Artefakts. Während interne Validatoren
Turing-spezifische Invarianten überprüfen, wird mit SYSIDE zusätzlich geprüft,
ob der erzeugte Modelltext durch eine von der Implementierung getrennte
SysML-v2-Toolchain verarbeitet werden kann. Ein erfolgreicher technischer
Validierungsschritt bestätigt damit die untersuchte technische
Kompatibilität, nicht jedoch automatisch die fachliche Richtigkeit des
zugrunde liegenden Engineering-Inhalts.

Fachliche Evidenz entsteht dagegen an den expliziten Human-Authority-Grenzen
des Systems. Dazu gehören der Human Engineering Review vor der Promotion zu
Approved Engineering Information, Human Model Placement, Target-Model
Formulation und Model Quality Review sowie der abschließende Final Human Model
Review.

Diese Reviews besitzen unterschiedliche Prüfgegenstände. Ein früher Human
Engineering Review bewertet die aus einer Source abgeleitete fachliche
Information. Model Placement und Target-Model Formulation adressieren dagegen
deren Überführung in eine geeignete Modellstruktur und Zielrepräsentation.
Der Final Human Model Review betrachtet schließlich das vollständige erzeugte
und technisch validierte Modell vor seiner Freigabe.

Als zusätzliche fachliche Evidenz wird eine externe Practitioner-Demonstration
eingesetzt. Dabei wird nicht versucht, aus einer kleinen Zahl beteiligter
Praktiker statistisch generalisierbare Aussagen abzuleiten. Ziel ist vielmehr
eine strukturierte Plausibilisierung, ob der demonstrierte Engineering
Information Flow, die Trennung zwischen KI-Vorschlag und Human Authority sowie
Provenance und Traceability für externe fachkundige Personen nachvollziehbar
sind. Ablauf und Beobachtungskriterien dieser Demonstration werden vor ihrer
Durchführung festgelegt und in Abschnitt~\ref{sec:safe_demo_verification}
beschrieben.

Die verschiedenen Evidenzformen ergänzen sich damit, ohne sich gegenseitig zu
ersetzen. Erst ihre Kombination ermöglicht es, sowohl die technische
Konsistenz der prototypischen Realisierung als auch die vorgesehenen
fachlichen Kontroll- und Freigabegrenzen der Architektur zu untersuchen.

% =============================================================================
\subsection{Automatisierte und formative Verifikation}
\label{sec:automated_formative_verification}
% =============================================================================

Die technische Verifikation des Proof of Concept erfolgte auf mehreren
aufeinander aufbauenden Ebenen. Einzelne Regeln und Datenverträge wurden
isoliert geprüft, während Integrations- und Workflowtests die Übergänge
zwischen mehreren Verarbeitungsschritten adressierten. Nach Änderungen an
einem begrenzten Systembereich wurden zunächst fokussierte Tests ausgeführt,
bevor eine vollständige Repository-Regression die bereits akzeptierten
Funktionalitäten erneut überprüfte.

Die Strukturierung der Prüfaktivitäten in isolierte Tests, Integrationstests
und übergreifende Regression folgt dem Grundprinzip, unterschiedliche
Testgegenstände und Testebenen getrennt zu untersuchen und ihre Ergebnisse
anschließend auf Systemebene zusammenzuführen. Die Planung und Durchführung
solcher gestufter Evaluationen ist ein etabliertes Vorgehen in der
empirischen Software-Engineering-Forschung \cite{Wohlin.2012}.

Diese automatisierte Testhierarchie wurde durch eine formative
Verifikationslogik ergänzt. Während kontrollierter End-to-End-Durchläufe
auftretende Abweichungen wurden als explizite Findings beziehungsweise
Blocker dokumentiert. Ein fehlgeschlagener Verarbeitungsschritt führte damit
nicht zu einer stillen Anpassung des Testablaufs, sondern zu einem
nachvollziehbaren Zyklus aus Befund, Korrektur und Retest.

Die automatisierte Verifikation bildet damit sowohl eine technische
Absicherung einzelner Contracts als auch eine Voraussetzung für die in den
folgenden Abschnitten beschriebenen End-to-End-Untersuchungen.


\subsubsection{Unit-, Contract- und Integrationsverifikation}
\label{sec:unit_contract_integration_tests}

Die kleinste Verifikationsebene bilden Unit- und Contract-Tests. Sie prüfen
isolierte technische Regeln und Datenstrukturen, beispielsweise die Bildung
und Validierung technischer Identitäten, die Konsistenz von Manifesten und
Fingerprints sowie die Bindung von Entscheidungen an konkrete
Artefaktzustände.

Von besonderer Bedeutung sind dabei fail-closed definierte Contracts.
Unvollständige, inkonsistente oder nicht mehr zum konsumierten
Artefaktzustand passende Informationen dürfen nicht implizit downstream
akzeptiert werden. Dies betrifft insbesondere Authority Boundaries, an denen
ein maschinell erzeugter oder zuvor autorisierter Zustand in die nächste
Verarbeitungsstufe überführt wird.

Repräsentative Testfamilien adressieren deshalb unter anderem Processing
Contexts und Artifact Identities, Human Decisions und Approved Engineering
Information, Project Fit und Model Candidate Sets, Model Placement und Model
Assembly sowie die nachgelagerte Modell- und Generierungskette. Die Tests
prüfen dabei nicht nur zulässige Zustände, sondern gezielt auch negative
Fälle, in denen fehlende Bindungen, ungültige Fingerprints oder unvollständige
Authority Sets eine Fortsetzung blockieren müssen.

Auf der nächsten Ebene überprüfen Integrations- und Workflowtests das
Zusammenspiel mehrerer solcher Contracts. Im Vordergrund stehen dabei die
Handoffs zwischen den in Kapitel~5 beschriebenen Verarbeitungsschritten.
Dazu gehören insbesondere

\begin{quote}
Processing und Evidence Formation\\
$\rightarrow$ Human Engineering Review\\
$\rightarrow$ Approved Engineering Information\\
$\rightarrow$ Project Fit\\
$\rightarrow$ Model Candidate Derivation und Human Placement\\
$\rightarrow$ Model Assembly Review und Base IEM\\
$\rightarrow$ Target-Model Formulation Authority\\
$\rightarrow$ Model-Quality Refinement und Human Model Quality Review\\
$\rightarrow$ Successor IEM\\
$\rightarrow$ SysML-v2-Generierung\\
$\rightarrow$ Generated Artifact Validation\\
$\rightarrow$ Final Human Model Review und Publication.
\end{quote}

Ein erfolgreicher Integrations- oder Workflowtest liefert dabei nur Evidenz
für den tatsächlich geprüften Ausschnitt. Insbesondere wird aus einem
erfolgreichen lokalen Handoff nicht auf einen erfolgreichen vollständigen
End-to-End-Durchlauf geschlossen. Diese übergreifende Evidenz wird im kontrollierten
WP-12-End-to-End-Dry-Run in Abschnitt~\ref{sec:wp12_e2e}
betrachtet.


\subsubsection{Focused und Complete Regression}
\label{sec:regression_verification}

Änderungen an Architekturcontracts oder Implementierungslogik wurden zunächst
durch eine \emph{Focused Regression} abgesichert. Die Begriffe
\emph{Focused Regression} und \emph{Complete Repository Regression}
bezeichnen dabei projektspezifische Ausprägungen der allgemeinen
Regressionstestlogik. Sie werden in dieser Arbeit nicht als eigenständige
normative Testkategorien verstanden. Dabei werden gezielt
diejenigen Testfamilien ausgeführt, die den korrigierten Verarbeitungsschritt
und seine unmittelbar angrenzenden Handoffs betreffen.

Dieser Schritt dient einer schnellen und lokal begrenzten Rückmeldung darüber,
ob die beabsichtigte Korrektur den betroffenen Contract wieder erfüllt. Ein
Focused Regression Pass stellt jedoch keine ausreichende Evidenz dafür dar,
dass unveränderte Bereiche des Systems durch die Korrektur unbeeinflusst
geblieben sind.

Nach akzeptierten Änderungen wurde deshalb zusätzlich eine
\emph{Complete Repository Regression} durchgeführt. Diese umfasst die
gesamte automatisierte Testsuite des Repositorys und prüft, ob bereits
akzeptierte technische Eigenschaften nach einer Änderung weiterhin bestehen.

Der finale für die Thesis akzeptierte Integrationsstand umfasst
6335 bestandene und 12 übersprungene automatisierte Tests. Diese Zahl
beschreibt den Umfang der abschließend ausgeführten Regressionsevidenz. Sie
wird nicht als Qualitätskennzahl des erzeugten Engineering-Modells
interpretiert. Insbesondere kann eine vollständige technische Regression weder
fachliche Engineering Correctness noch die Plausibilität eines erzeugten
SysML-v2-Modells nachweisen.

Focused und Complete Regression erfüllen damit unterschiedliche Funktionen.
Die erste lokalisiert die technische Evidenz einer konkreten Änderung, die
zweite schützt den bereits erreichten integrierten Systemzustand vor
unbeabsichtigten Regressionen.


\subsubsection{Findings, Blocker und Retest-Logik}
\label{sec:finding_retest_logic}

Die automatisierten Tests wurden während der formativen Systemverifikation
durch eine explizite Finding- und Blocker-Logik ergänzt. Dadurch konnten auch
Fehlerzustände dokumentiert werden, die erst im Zusammenspiel mehrerer
Verarbeitungsschritte beziehungsweise während eines kontrollierten
End-to-End-Durchlaufs sichtbar wurden.

Für die durchgeführten Verifikationsläufe wurden die Ergebniszustände
\emph{PASS}, \emph{PASS WITH FINDINGS} und
\emph{FAILED WITH BLOCKER} unterschieden. Ein nicht blockierendes Finding
konnte dokumentiert bleiben, ohne den gesamten Prüfpfad ungültig zu machen.
Ein Blocker verhinderte dagegen die Fortsetzung beziehungsweise Akzeptanz des
betroffenen Gates, bis seine Ursache adressiert und der relevante
Verarbeitungsschritt erneut geprüft worden war.

Die konkrete Anwendung dieser Ergebnis- und Retest-Logik wurde für WP-12
vor der Durchführung in einem detaillierten Testprotokoll festgelegt. Das
vollständige Protokoll ist in Anhang~\ref{app:wp12_test_protocol}
wiedergegeben.

Die zugrunde liegende Retest-Logik folgt dabei grundsätzlich der Abfolge

\begin{quote}
kontrollierter Prüfzustand\\
$\rightarrow$ Finding oder Blocker\\
$\rightarrow$ Ursachenanalyse\\
$\rightarrow$ begrenzte Korrektur\\
$\rightarrow$ Focused Regression\\
$\rightarrow$ Retest des betroffenen Gates\\
$\rightarrow$ Complete Repository Regression.
\end{quote}

Ein wesentliches Merkmal dieses Vorgehens ist der Erhalt der ursprünglichen
Fehlerevidenz. Ein blockierter Projekt- oder Verarbeitungszustand wurde nicht
durch einen neu erzeugten, scheinbar fehlerfreien Ausgangszustand ersetzt.
Stattdessen blieben Project, Run, relevante Artefakte und dokumentierter Befund
erhalten und bildeten die Ausgangsbasis für die nachfolgende Korrektur und den
Retest.

Die Findings besitzen dabei unterschiedliche technische Tragweiten. Ein
Befund kann auf einen lokalen Implementierungsfehler zurückzuführen sein,
aber auch auf einen zu permissiven Contract, eine unzureichende
Artefaktidentität, eine fehlende Processing Stage oder eine falsch
abgegrenzte Authority Responsibility hinweisen. Nur wenn ein Befund die
Architektur- oder Responsibility-Grenzen berührt, ist daraus eine
architekturrelevante Korrektur beziehungsweise Design-Rationale abzuleiten.

Die konkreten Befunde und ihre beobachteten Auswirkungen werden an dieser
Stelle noch nicht als Ergebnisse zusammengefasst. Die für die
Single-Source- und Multi-Source-Verifikation relevanten Fälle werden in den
folgenden Abschnitten anhand der jeweiligen Versuchsanordnung eingeordnet.
Die aggregierte Ergebnisdarstellung erfolgt anschließend in Kapitel~7.

\subsection{Kontrollierte WP-12-End-to-End-Verifikation}
\label{sec:wp12_e2e}

Die in den vorangegangenen Abschnitten beschriebenen automatisierten Tests
prüfen einzelne Contracts, Handoffs und Verarbeitungskomponenten. Ergänzend
wurde mit WP-12 ein kontrollierter End-to-End-Dry-Run definiert, der den
zusammenhängenden Engineering Information Flow unter Multi-Source-Bedingungen
untersucht.

WP-12 wurde von Beginn an als Multi-Document-Szenario konzipiert. Mehrere
getrennte Legacy-Dokumente werden innerhalb eines gemeinsamen Project-Kontexts
registriert, verarbeitet und menschlich bewertet. Anschließend wird geprüft,
ob die daraus autorisierten Engineering-Informationszustände gemeinsam in den
projektweiten Modellableitungspfad eingehen können.

Ziel des Dry Runs ist damit nicht allein der Nachweis, dass einzelne
Verarbeitungsschritte technisch ausführbar sind. Untersucht wird, ob die in
Kapitel~5 beschriebene Verarbeitungskette unter kontrollierten Bedingungen
durchgängig durchlaufen werden kann und dabei insbesondere Source Provenance,
Human Authority und die Trennung verschiedener Informationszustände erhalten
bleiben.

Das Testdesign, die verwendeten Eingangsdaten und die erwarteten
Engineering-Eigenschaften wurden vor der Durchführung festgelegt. Die
vollständigen Testdaten sind in Anhang~\ref{app:vv_test_data} dokumentiert.
Das detaillierte WP-12-Testprotokoll sowie der zugehörige formative
Auswertungsbericht sind in Anhang~\ref{app:wp12_test_protocol} enthalten.


\subsubsection{Testdatenbasis und Szenariokonstruktion}
\label{sec:wp12_test_data}

Als kontrollierte Testdatenbasis wurde ein synthetischer
Legacy-Dokumentensatz zum Anwendungsszenario einer
Remote-Microscope-Collaboration erstellt. Der Datensatz besteht aus vier
getrennten Quelldokumenten, die unterschiedliche Informationsperspektiven,
Formalisierungsgrade und Engineering-Inhalte repräsentieren.

Die erste Source beschreibt den Produktkontext und die beabsichtigte
Remote-Collaboration auf informeller Ebene. Neben bereits formulierten
Produkteigenschaften enthält sie bewusst noch offene Fragestellungen, unter
anderem zur Aufbewahrungsdauer von Session-Informationen und zu akzeptablen
Grenzen der Verbindungsqualität.

Die zweite Source beschreibt den Nutzerworkflow zwischen Microscope Operator
und Remote Expert. Sie enthält sowohl den regulären Ablauf eines explizit
autorisierten Control Transfers als auch ein davon abweichendes
Recovery-Verhalten bei einem Verbindungsverlust des Remote Expert.

Die dritte Source enthält einen Satz expliziter, jedoch als Entwurf
gekennzeichneter System Requirements. Darin werden unter anderem ein einzelner
aktiver Controller, ein expliziter normaler Control Transfer,
Controller-Visibility, die Revocation von Remote Authority bei
Verbindungsverlust sowie die Traceability von Session- und
Control-Authority-Änderungen adressiert.

Die vierte Source beschreibt einen informellen technischen Architekturstand.
Sie benennt Verantwortlichkeiten für Streaming, Control und
Session Traceability, legt deren konkrete Deployment-Zuordnung jedoch bewusst
nicht fest.

Der Datensatz bildet damit keine bereits konsolidierte oder fachlich
vollständige Systemspezifikation. Seine Dokumente enthalten bewusst
überlappende Aussagen, komplementäre Informationen, unterschiedliche
Abstraktionsebenen, offene Fragen und noch nicht abschließend getroffene
Engineering-Entscheidungen.

Ein Beispiel für dokumentübergreifende Überlappung ist die Remote-Viewing-
Funktionalität, die aus unterschiedlichen Perspektiven in mehreren Sources
beschrieben wird. Andere Engineering-Inhalte verteilen sich komplementär über
mehrere Dokumente. Die Control-Funktionalität wird beispielsweise auf Produkt-,
Workflow-, Requirement- und technischer Ebene beschrieben, ohne dass eine
einzelne Source den gesamten Sachverhalt enthält.

Der Datensatz enthält darüber hinaus gezielte semantische Spannungen. Der
reguläre Control Transfer erfordert eine explizite Entscheidung des
Microscope Operator, während bei einem Verbindungsverlust eine automatische
Rücknahme der Remote Authority beschrieben wird. Beide Aussagen sind fachlich
miteinander vereinbar, dürfen jedoch nicht zu einer identischen Regel
normalisiert oder fälschlich als Widerspruch interpretiert werden.

Zusätzlich wurden bestimmte Informationen bewusst nicht spezifiziert. Dazu
gehören quantitative Vorgaben zu Latenz und Bandbreite, eine konkrete
Retention Time, die Deployment-Topologie, verwendete Netzwerkprotokolle,
Persistenztechnologien und Cybersecurity-Zielwerte. Dadurch kann geprüft
werden, ob fehlende Engineering-Information als solche erhalten bleibt und
nicht durch unbelegte Angaben ergänzt wird.

Die Heterogenität des verwendeten Testdatensatzes bezieht sich damit auf
Informationsart, Engineering-Perspektive, Formalisierungsgrad und
Informationsgehalt. Technisch liegen alle vier Sources als Markdown-Dateien
vor. Die Untersuchung adressiert daher keine allgemeine Verarbeitung
heterogener Dateiformate wie PDF-, Office- oder proprietärer
Engineering-Formate.

Die vollständigen und tatsächlich verwendeten Quelldokumente werden in
Anhang~\ref{app:vv_test_data} wiedergegeben. Dadurch bleibt die
Eingangsdatenbasis der V\&V unabhängig vom generierten Ergebnis vollständig
nachvollziehbar.


\subsubsection{Testdesign und Versuchsdurchführung}
\label{sec:wp12_test_design}

Der WP-12-Dry-Run wurde als kontrollierter, schrittweise geführter
End-to-End-Durchlauf angelegt. Die Versuchsdurchführung kombiniert den
Human-facing Workflow des Demonstrators mit automatisierter technischer
Verifikation und expliziten Human-Authority-Entscheidungen.

Zu Beginn werden die vier Testdokumente getrennt innerhalb desselben
Project-Kontexts registriert. Für jede Source muss eine eigenständige
Source Identity bestehen bleiben. Die Registrierung oder Verarbeitung einer
weiteren Source darf bereits vorhandene Source-Zustände nicht ersetzen oder
überschreiben.

Anschließend werden die Sources jeweils über den regulären Processing-Pfad
verarbeitet. Die entstehenden Evidence- und Interpretation-Zustände werden
nicht automatisch zu autorisierter Engineering Information. Stattdessen
werden sie dem vorgesehenen Human Engineering Review zugeführt. Erst
explizit menschlich akzeptierte Informationen dürfen in einen autorisierten
Downstream-Zustand überführt werden.

Nach Abschluss der source-lokalen Reviews wird geprüft, ob die autorisierten
Informationszustände mehrerer Sources innerhalb desselben Projects gemeinsam
für die nachgelagerte Modellableitung verfügbar gemacht werden können. Dabei
müssen ihre jeweiligen Source- und Authority-Bindungen erhalten bleiben.

Der weitere Testpfad folgt der in Kapitel~5 beschriebenen Modellableitung über
Project Fit, Model Candidate Derivation und Human Model Placement. Daran
schließen sich Model Assembly Review, die Bildung des Base Internal
Engineering Model, Target-Model Formulation, Model-Quality Refinement und
Human Model Quality Review an. Das daraus autorisierte Successor Internal
Engineering Model bildet den Eingang für die deterministische
SysML-v2-Generierung.

Der vollständige geprüfte Informationsfluss lässt sich damit in
vereinfachter Form als

\begin{quote}
mehrere Engineering Sources\\
$\rightarrow$ source-lokales Processing\\
$\rightarrow$ Human Engineering Review\\
$\rightarrow$ autorisierte source-lokale Engineering Information\\
$\rightarrow$ Project Fit\\
$\rightarrow$ projektweite Model Candidate Derivation\\
$\rightarrow$ Human Model Placement\\
$\rightarrow$ Model Assembly Review und Base IEM\\
$\rightarrow$ Target-Model Formulation\\
$\rightarrow$ Model-Quality Refinement und Human Model Quality Review\\
$\rightarrow$ Human-authorized Successor IEM\\
$\rightarrow$ deterministische SysML-v2-Generierung\\
$\rightarrow$ Generated Artifact Validation\\
$\rightarrow$ Final Human Model Review\\
$\rightarrow$ kontrollierte Publication
\end{quote}

zusammenfassen.

Für die einzelnen Stufen definiert das Testprotokoll konkrete
Preconditions, erwartete Zustände und Akzeptanzbedingungen. Ein technischer
oder fachlicher Gate-Zustand wird nur dann fortgesetzt, wenn die für den
jeweiligen Testschritt erforderlichen Bedingungen erfüllt sind. Auftretende
Abweichungen werden nach der in Abschnitt~\ref{sec:finding_retest_logic}
beschriebenen Finding- und Blocker-Logik behandelt.

Der vollständige Ablauf mit den einzelnen Testfällen, erwarteten Zuständen und
Dokumentationsfeldern wird im Anhang
über das originale WP-12-Testprotokoll bereitgestellt. Der Haupttext
beschränkt sich auf die Versuchsanordnung und diejenigen Prüfbedingungen,
die für die spätere Einordnung der Ergebnisse erforderlich sind.


\subsubsection{Semantische und architektonische Akzeptanzkriterien}
\label{sec:wp12_acceptance_criteria}

Zusätzlich zu den technischen Testbedingungen wurde für WP-12 ein
\emph{Expected Engineering Contract} definiert. Dieser legt keine exakte
Zielstruktur des erzeugten SysML-v2-Modells fest. Bewertet wird vielmehr, ob
die aus den Testdaten erwartbare Engineering-Bedeutung, ihre Provenance und
die vorgesehenen Authority Boundaries im Verarbeitungsergebnis erhalten
bleiben.

Auf inhaltlicher Ebene wird erwartet, dass die zentralen Rollen und
Engineering-Sachverhalte des Testdatensatzes im resultierenden
Informationszustand erhalten bleiben. Dazu gehören insbesondere Microscope
Operator und Remote Expert sowie die Kollaborationssession, Remote Viewing,
Control Request und Control Transfer, Remote Microscope Control,
Connection-Loss-Recovery und Session Traceability.

Gleichzeitig müssen wesentliche Constraints der Quelldokumente erhalten
bleiben. Dazu zählen insbesondere die Begrenzung auf einen aktiven Controller,
die explizite Entscheidung beim regulären Control Transfer, die Sichtbarkeit
des aktuellen Controllers, die Rücknahme der Remote Authority bei
Verbindungsverlust und die Aufzeichnung relevanter Session- und
Control-Authority-Änderungen.

Für die dokumentübergreifende Verarbeitung werden darüber hinaus mehrere
spezifische Kriterien geprüft. Semantisch überlappende Aussagen sollen als
inhaltlich zusammengehörig erkannt werden können, ohne dass ihre jeweilige
Source Provenance verloren geht. Komplementäre Informationen aus
unterschiedlichen Sources müssen gemeinsam zur Modellableitung beitragen
können, auch wenn keine einzelne Source den vollständigen Engineering-
Sachverhalt enthält.

Semantische Unterschiede dürfen dagegen nicht allein zugunsten einer
automatischen Vereinheitlichung entfernt werden. Insbesondere müssen der
reguläre menschlich autorisierte Control Transfer und das automatische
Recovery-Verhalten bei Verbindungsverlust unterscheidbar bleiben. Verbleibende
Mehrdeutigkeit ist als Review-Gegenstand zu behandeln und darf nicht
stillschweigend durch die automatisierte Verarbeitung entschieden werden.

Ein weiteres Akzeptanzkriterium ist der Erhalt source-spezifischer
Information. Die gemeinsame Verarbeitung mehrerer Sources darf nicht dazu
führen, dass eine Aussage allein deshalb verloren geht, weil sie nur in einem
Dokument vorkommt. Ebenso darf aus der Reihenfolge der Registrierung oder
Verarbeitung keine implizite fachliche Priorisierung einer Source entstehen.

Die bewusst fehlenden Informationen des Testdatensatzes bilden
Negativkriterien. Nicht belegte quantitative Werte oder technische
Festlegungen dürfen nicht als autorisierte Engineering-Aussagen erzeugt
werden. Fehlende Information darf explizit kenntlich gemacht werden, ohne die
Verarbeitung der tatsächlich belegten Engineering-Inhalte zu verhindern.

Schließlich wird die Multi-Source-Fähigkeit nicht bereits dadurch als
nachgewiesen betrachtet, dass mehrere Dokumente technisch eingelesen werden
können. Erforderlich ist vielmehr, dass mehrere getrennte und
provenance-gebundene Source-Zustände unter Erhalt der vorgesehenen
Human-Authority-Grenzen in einen gemeinsamen projektweiten
Modellableitungspfad eingehen können.

Die konkreten während WP-12 beobachteten Ergebnisse werden erst in
Kapitel~7 dargestellt. Dort werden sowohl erfolgreiche Prüfungen als auch
während der Durchführung aufgetretene Abweichungen zusammengeführt. Die
architektonische Bedeutung dieser Beobachtungen und daraus entstandene
Anpassungen werden anschließend in Kapitel~8 diskutiert.

\subsection{SysML-v2-Validierung und Human Release}
\label{sec:sysml_validation_release}

Der WP-12-End-to-End-Pfad endet nicht mit der Erzeugung eines
SysML-v2-Artefakts. Vor einer abschließenden Publication wird das generierte
Artefakt zunächst technisch validiert und anschließend einem eigenständigen
Final Human Model Review unterzogen.

Diese beiden Prüfungen adressieren unterschiedliche Fragestellungen. Die
technische Validierung untersucht, ob das erzeugte Artefakt die für die
Generierung und Zielnotation festgelegten technischen Bedingungen erfüllt.
Der Final Human Model Review bewertet dagegen das vollständige resultierende
Modell als Engineering-Artefakt und bildet die Grundlage für die explizite
menschliche Freigabe zur Publication.

Die Trennung ist wesentlich für die in Kapitel~4 definierte
Authority-Struktur. Ein technisch valides Modell ist nicht automatisch
fachlich freigegeben. Umgekehrt kann eine menschliche Bewertung technische
Syntax-, Integritäts- oder Traceability-Prüfungen nicht ersetzen.


\subsubsection{Interne und externe Generated-Artifact-Validierung}
\label{sec:generated_artifact_validation}

Die technische Validierung wird auf dem vollständig generierten
SysML-v2-Artefakt durchgeführt. Prüfgegenstand ist damit nicht erneut der
semantische Engineering-Inhalt der vorgelagerten Verarbeitung, sondern die
konkrete aus dem autorisierten Internal Engineering Model erzeugte
Zielrepräsentation.

Die Validierung besteht aus zwei voneinander getrennten Ebenen. Zunächst
werden innerhalb des Turing Generators deterministische
Generated-Artifact-Checks durchgeführt. Sie prüfen Turing-spezifische
Invarianten wie Artifact Integrity, Fingerprint-Bindungen, verwendete
Generation Profiles, zulässige Target-Notation-Konstrukte, strukturelle
Konsistenz, Relationship Endpoints und die Vollständigkeit der
Traceability-Informationen.

Diese interne Prüfung ist bewusst auf den durch den Proof of Concept
unterstützten Generierungsumfang begrenzt. Sie implementiert keine
vollständige SysML-v2-Sprachverarbeitung. Die normative Grundlage der
verwendeten Zielnotation bildet die SysML-v2-Sprachspezifikation
\cite{SysMLv2}.

Ergänzend wird das vollständige generierte Artefakt durch eine von der
Turing-Implementierung getrennte Toolchain geprüft. Hierfür wird der
SYSIDE Modeler über eine kontrollierte externe Validator-Schnittstelle
verwendet. Dadurch lässt sich untersuchen, ob die tatsächlich erzeugte
SysML-v2-Repräsentation auch außerhalb der generatorinternen
Validierungslogik verarbeitet werden kann.

Die beiden Validierungsebenen ersetzen sich nicht gegenseitig. Interne
Validatoren können projektspezifische Contracts, Traceability und
Generation-Profile überprüfen, die einem externen SysML-v2-Werkzeug nicht
bekannt sind. Die externe Validierung liefert dagegen unabhängige Evidenz
für die Parser- und Toolkompatibilität des erzeugten Artefakts.

Die technische Validierung arbeitet beobachtend. Ein als fehlerhaft erkanntes
Artefakt wird innerhalb dieses Verarbeitungsschritts nicht automatisch
verändert oder repariert. Änderungen müssen an der jeweils verantwortlichen
vorgelagerten Processing-, Model- oder Generation-Stufe vorgenommen und
anschließend erneut generiert und validiert werden.

Für das Validierungsergebnis werden die drei Zustände
\emph{valid}, \emph{invalid} und \emph{incomplete} unterschieden.
\emph{Invalid} bezeichnet einen Zustand, in dem die Validierung einen
blockierenden Defekt des Artefakts festgestellt hat. \emph{Incomplete}
wird verwendet, wenn eine erforderliche Validierungsaktivität nicht
vollständig ausgeführt werden konnte, beispielsweise weil die externe
Validator-Infrastruktur nicht verfügbar ist.

Beide Zustände bleiben fail-closed. Nur ein vollständig ausgeführter
Validierungspfad ohne blockierende Findings kann die technische
Publication Eligibility herstellen. Insbesondere darf die
Nichtverfügbarkeit eines externen Validators nicht als erfolgreicher
Validierungszustand interpretiert werden.

Ein erfolgreicher SYSIDE-Durchlauf besitzt zugleich eine klare
Claim Boundary. Er liefert Evidenz dafür, dass das erzeugte Artefakt durch
die verwendete externe SysML-v2-Toolchain verarbeitet werden kann. Daraus
folgt weder die fachliche Korrektheit des Engineering-Modells noch dessen
Eignung für einen beliebigen Engineering-Kontext.


\subsubsection{Final Human Model Review und kontrollierte Publication}
\label{sec:final_human_review_publication}

Nach der technischen Validierung wird das vollständig generierte Modell einem
eigenständigen Final Human Model Review zugeführt. Dieser Review unterscheidet
sich von den früheren Human-Authority-Grenzen des Informationsflusses.

Der Human Engineering Review bewertet source-grounded Engineering Information,
während Human Model Placement und Model Quality Review Entscheidungen über
deren Überführung in den autorisierten internen Modellzustand adressieren.
Der Final Human Model Review betrachtet dagegen erstmals das vollständige
resultierende Modell in seiner generierten Zielrepräsentation.

Für die Bewertung werden das erzeugte SysML-v2-Artefakt, der zugehörige
Validierungszustand und die vorhandene Traceability gemeinsam bereitgestellt.
Dadurch kann die fachliche Bewertung auf genau den Artefaktzustand bezogen
werden, der zuvor technisch validiert wurde.

Technisch ungültige oder unvollständig validierte Artefakte können weiterhin
als Review-Evidenz betrachtet werden, beispielsweise um eine erforderliche
Korrektur zu identifizieren. Sie sind jedoch nicht zur abschließenden
Publication berechtigt.

Eine materielle Änderung am Modell erzeugt einen neuen Artefaktzustand.
Dadurch werden auch eine erneute Generation, eine erneute technische
Validierung und eine neue Final-Review-Revision erforderlich. Frühere
Validierungs- oder Freigabeentscheidungen werden nicht auf veränderte
Artefakte übertragen.

Die abschließende Human Release Decision ist eine eigenständige
Authority-Entscheidung. Sie darf weder aus einem erfolgreichen
Validierungsergebnis noch aus der Abwesenheit technischer Findings,
KI-Konfidenz oder Agreement zwischen automatisierten Interpretationen
abgeleitet werden.

Eine Publication ist erst zulässig, wenn der zu veröffentlichende
Artefaktzustand technisch valide ist, der zugehörige Validierungszustand
exakt zu diesem Artefakt gehört und eine explizite menschliche Freigabe für
genau diesen Review-Zustand vorliegt.

Der Abschluss des V\&V-Pfads folgt damit der Sequenz

\begin{quote}
Generated SysML-v2 Artifact\\
$\rightarrow$ interne Generated-Artifact-Validierung\\
$\rightarrow$ externe SYSIDE-Validierung\\
$\rightarrow$ Final Human Model Review\\
$\rightarrow$ explizite Human Release Decision\\
$\rightarrow$ kontrollierte Publication.
\end{quote}

Der veröffentlichte Output wird anschließend als versionierter und
unveränderlicher Projektzustand behandelt. Generierte, validierte und
veröffentlichte Artefakte bleiben damit begrifflich und technisch
unterschiedliche Zustände.

Die in diesem Abschnitt beschriebenen Prüfschritte definieren die
Akzeptanzlogik des abschließenden Modellpfads. Welche Zustände während der
durchgeführten V\&V tatsächlich erreicht wurden, wird erst in Kapitel~7
dargestellt.

\subsection{Externe Stakeholder-Demonstration}
\label{sec:safe_demo_verification}

Ergänzend zur technischen und kontrollierten End-to-End-Verifikation wird der
Proof of Concept in einer externen Stakeholder-Demonstration vorgestellt.
Ziel ist es, die Nachvollziehbarkeit des realisierten Engineering Information
Flow sowie der vorgesehenen Human-Authority- und Traceability-Mechanismen aus
Sicht fachkundiger externer Stakeholder zu untersuchen.

Die Demonstration stellt dabei keinen kontrollierten experimentellen Versuch
und keine statistische Usability-Evaluation dar. Sie dient vielmehr als
ergänzende fachliche Plausibilisierung des entwickelten Ansatzes. Um eine
nachträgliche Anpassung der Bewertungskriterien an beobachtete Ergebnisse zu
vermeiden, werden Ablauf, Systembaseline, Beobachtungskriterien und
Feedbackfragen vor der Durchführung festgelegt.


\subsubsection{Zielsetzung und Einordnung}
\label{sec:stakeholder_demo_goal}

Die Stakeholder-Demonstration adressiert insbesondere die Frage, ob der in
Kapitel~5 beschriebene Informationsverarbeitungspfad für externe fachkundige
Personen nachvollziehbar bleibt, obwohl mehrere automatisierte, semantische und
menschlich autorisierte Verarbeitungsschritte kombiniert werden.

Im Mittelpunkt stehen deshalb nicht einzelne Implementierungsdetails, sondern
die übergreifenden Engineering-Prinzipien des Proof of Concept. Dazu zählen die
Trennung zwischen KI-generiertem Vorschlag und Human Authority, die
Nachvollziehbarkeit der Herkunft einer Engineering-Aussage, die Verarbeitung
mehrerer Sources innerhalb eines gemeinsamen Project-Kontexts sowie die
Überführung autorisierter Engineering Information in ein technisch
validierbares SysML-v2-Modell.

Die Demonstration dient zugleich als externe Plausibilisierung der
Human-facing Workflow-Gestaltung. Dabei wird untersucht, ob für die
teilnehmenden Stakeholder erkennbar bleibt, welche Information lediglich
automatisch vorgeschlagen wurde, welche Entscheidung bereits menschlich
autorisiert ist und an welchen Stellen weitere Engineering-Entscheidungen
erforderlich sind.

Aus der Demonstration werden keine statistisch generalisierbaren Aussagen über
Usability, Akzeptanz oder Engineering Correctness abgeleitet. Die erhobenen
Beobachtungen ergänzen vielmehr die technische V\&V um eine externe fachliche
Perspektive auf den demonstrierten Workflow.


\subsubsection{Systembaseline und Demonstrationsszenario}
\label{sec:stakeholder_demo_baseline}

Für die Demonstration wird ein vorab festgelegter Systemzustand verwendet.
Damit soll verhindert werden, dass während der Durchführung spontan ein
technisch oder inhaltlich günstigerer Ausgangszustand gewählt wird.

Vor der Demonstration werden mindestens die verwendete
Implementierungsbaseline, der Stand des Architekturmodells, die eingesetzten
Test-Sources und der vorbereitete Project-Zustand dokumentiert. Sofern für die
Reproduzierbarkeit relevant, werden zusätzlich der verwendete LLM-Provider und
das eingesetzte Modell festgehalten.

Die Demonstration verwendet denselben grundlegenden Engineering Information
Flow, der bereits Gegenstand der kontrollierten WP-12-Verifikation ist. Die
teilnehmenden Stakeholder erhalten damit keinen speziell für die Demonstration
vereinfachten fachlichen Verarbeitungspfad. Die Präsentation kann jedoch
bestehende persistierte Zustände verwenden, um die einzelnen
Verarbeitungsschritte in einer zeitlich geeigneten Form zeigen zu können.

Der geplante Demonstrationsablauf umfasst in vereinfachter Form

\begin{quote}
Source Registration\\
$\rightarrow$ Processing und source-grounded Evidence\\
$\rightarrow$ semantische Interpretation\\
$\rightarrow$ Human Engineering Review\\
$\rightarrow$ autorisierte Engineering Information\\
$\rightarrow$ Multi-Source- und Project-Level-Verarbeitung\\
$\rightarrow$ Model Candidate Derivation und Human Model Placement\\
$\rightarrow$ Internal Engineering Model\\
$\rightarrow$ Target-Model- und Model-Quality-Verarbeitung\\
$\rightarrow$ deterministische SysML-v2-Generierung\\
$\rightarrow$ technische Validierung\\
$\rightarrow$ Final Human Model Review und Publication.
\end{quote}

Abweichungen vom geplanten Ablauf werden während der Durchführung dokumentiert.
Der Ablauf wird nicht nachträglich verändert, um beobachtete Ergebnisse als
erfolgreicher erscheinen zu lassen.

% NACH DER DEMONSTRATION ERGÄNZEN:
% Datum / Ort:
% Teilnehmerzahl:
% Rollen bzw. fachlicher Hintergrund:
% Git Commit / Implementierungsbaseline:
% CATIA Baseline:
% verwendeter Project-Zustand:
% verwendete Sources:
% LLM Provider / Modell, sofern relevant:
% dokumentierte Abweichungen vom geplanten Ablauf:


\subsubsection{Beobachtungskriterien und Feedbackerhebung}
\label{sec:stakeholder_demo_criteria}

Die Beobachtungskriterien werden vor der Durchführung festgelegt. Dadurch
bleibt nachvollziehbar, welche Eigenschaften des Demonstrationsworkflows
gezielt betrachtet wurden und welche Rückmeldungen erst während der
Demonstration entstanden sind.

Für die Demonstration werden folgende Kriterien verwendet:

\begin{enumerate}
    \item \textbf{K1: Nachvollziehbarkeit des Gesamtworkflows.}
    Die wesentlichen Verarbeitungsschritte und ihre Abfolge sollen für die
    Stakeholder fachlich nachvollziehbar sein.

    \item \textbf{K2: Trennung von KI-Vorschlag und Human Authority.}
    Es soll erkennbar sein, welche Inhalte automatisch vorgeschlagen werden
    und welche Zustände erst durch eine explizite menschliche Entscheidung
    Engineering Authority erhalten.

    \item \textbf{K3: Nachvollziehbarkeit von Provenance und Traceability.}
    Die Herkunft einer Engineering-Aussage und ihre Verbindung zu
    nachgelagerten Modellartefakten sollen nachvollziehbar sein.

    \item \textbf{K4: Nachvollziehbarkeit der Multi-Source-Verarbeitung.}
    Es soll verständlich sein, wie mehrere Sources zu einem gemeinsamen
    Project beitragen, ohne ihre jeweilige Herkunft zu verlieren.

    \item \textbf{K5: Plausibilität der Human-Review-Grenzen.}
    Die vorgesehenen menschlichen Review- und Freigabeschritte sollen aus
    Engineering-Sicht an nachvollziehbaren Stellen des Informationsflusses
    liegen.

    \item \textbf{K6: Fachliche Plausibilität des erzeugten
    SysML-v2-Ergebnisses.}
    Das demonstrierte Modellergebnis soll im Kontext der gezeigten
    Eingangsinformation fachlich nachvollziehbar sein.

    \item \textbf{K7: Wahrgenommene Risiken und fehlende Informationen.}
    Die Stakeholder sollen auf Stellen hinweisen können, an denen zusätzliche
    Information, weitere Kontrolle oder eine andere Form der Absicherung für
    einen realen Engineering-Einsatz erforderlich wäre.
\end{enumerate}

Neben den Beobachtungskriterien werden wenige offene Feedbackfragen verwendet.
Sie sollen zusätzliche fachliche Einschätzungen ermöglichen, ohne eine
bestimmte Antwort vorzugeben. Vorgesehen sind insbesondere folgende Fragen:

\begin{quote}
Welche Stelle des Workflows wirkt aus Ihrer Sicht fachlich am wenigsten
belastbar?

An welchen Stellen ist aus Ihrer Sicht zwingend eine menschliche Entscheidung
erforderlich?

Welche zusätzliche Information wäre für eine reale Engineering-Freigabe
erforderlich?

Ist die Herkunft einer erzeugten Modellaussage für Sie ausreichend
nachvollziehbar?

Welche Risiken oder Schwachstellen sehen Sie bei einer Übertragung des Ansatzes
auf reale Engineering-Projekte?
\end{quote}

Die Rückmeldungen werden unabhängig davon dokumentiert, ob sie den
Architekturansatz bestätigen oder auf Schwachstellen hinweisen. Positive und
kritische Beobachtungen werden nach denselben Kriterien erfasst.


\subsubsection{Dokumentation und Claim Boundary}
\label{sec:stakeholder_demo_claim_boundary}

Der geplante und der tatsächlich durchgeführte Ablauf werden getrennt
dokumentiert. Nach der Demonstration werden insbesondere Abweichungen vom
vorgesehenen Szenario, technische Auffälligkeiten sowie die Rückmeldungen zu
den vorab definierten Kriterien festgehalten.

Die Stakeholder-Demonstration kann dadurch Evidenz dafür liefern, ob
fachkundige externe Personen den demonstrierten Informationsfluss, die
Human-Authority-Grenzen und die Traceability-Mechanismen nachvollziehen und
fachlich bewerten können.

Der daraus ableitbare Claim bleibt jedoch begrenzt. Die Demonstration stellt
keinen Nachweis statistischer Usability, allgemeiner Nutzerakzeptanz,
universeller Engineering Correctness oder Production Readiness dar. Ebenso
kann aus einer positiven Bewertung einzelner Stakeholder nicht auf die
Eignung des Ansatzes für beliebige Engineering-Domänen geschlossen werden.

Die beobachteten Rückmeldungen werden in Kapitel~7 als Ergebnisse der
Stakeholder-Demonstration dargestellt. Ihre Bedeutung für die
Informationsverarbeitungsarchitektur, die Grenzen der Untersuchung und die
Forschungsfragen wird anschließend in Kapitel~8 diskutiert.

\subsection{Zuordnung der V\&V-Evidenz zu den Forschungsfragen}
\label{sec:vv_rq_mapping}

Die in diesem Kapitel beschriebenen Verifikations- und
Validierungsaktivitäten dienen nicht jeweils isoliert der Beantwortung einer
einzelnen Forschungsfrage. Vielmehr tragen unterschiedliche Evidenzformen zu
verschiedenen Aspekten der Haupt- und Teilforschungsfragen bei. Die
Forschungsfragen selbst werden an dieser Stelle noch nicht beantwortet. Die
folgende Zuordnung legt lediglich fest, welche V\&V-Evidenz in Kapitel~8 für
die spätere Beantwortung herangezogen wird.

Für die Hauptforschungsfrage ist insbesondere die Kombination der
unterschiedlichen Evidenzebenen relevant. Die automatisierte Verifikation
adressiert die Einhaltung zentraler Contracts und Authority Boundaries. Der
kontrollierte WP-12-End-to-End-Dry-Run untersucht den zusammenhängenden
Engineering Information Flow unter Multi-Source-Bedingungen. Die
Generated-Artifact-Validierung und die externe SYSIDE-Prüfung liefern
zusätzliche Evidenz für die technische Konsistenz des erzeugten
SysML-v2-Artefakts. Human Review und die externe Stakeholder-Demonstration
ergänzen diese technische Evidenz um die fachliche Perspektive.

Für TF1 sind vor allem diejenigen Prüfungen relevant, die den Umgang mit
Eingangsdaten und deren Eignung für die nachgelagerte Verarbeitung
adressieren. Dazu gehören Source Registration, Source Preparation,
source-grounded Evidence, die fail-closed Behandlung unvollständiger oder
nicht belegter Information sowie Project Fit. Die gezielt heterogen
konstruierten WP-12-Testdaten ermöglichen zusätzlich die Untersuchung von
unterschiedlichen Informationsarten, Formalisierungsgraden, semantischen
Überlappungen und bewusst fehlenden Engineering-Informationen.

TF2 wird insbesondere durch die Evidenz an den verschiedenen
Human-Authority-Grenzen adressiert. Dazu zählen Human Engineering Review,
Human Model Placement, Target-Model Formulation, Human Model Quality Review
und Final Human Model Review. Ergänzend untersucht die
Stakeholder-Demonstration, ob die Trennung zwischen automatisiertem Vorschlag
und menschlicher Entscheidung auch für externe fachkundige Personen
nachvollziehbar bleibt.

Für TF3 sind vor allem diejenigen Prüfungen relevant, die kontrollierte
semantische und strukturelle Leitplanken adressieren. Dazu gehören die
source-grounded Verarbeitung, die Project-Fit-Bedingungen, die
profilgebundene Modellableitung, die Trennung zwischen semantischer
Interpretation und deterministischer Modellverarbeitung sowie die
Validierung des generierten SysML-v2-Artefakts gegen die festgelegte
Target Notation und die zugehörigen Generation- und Validation Contracts.

Die Zuordnung macht zugleich sichtbar, dass keine einzelne Evidenzform für
sich einen vollständigen Nachweis der Forschungsfragen liefert. Automatisierte
Tests können technische Contract-Erfüllung zeigen, aber keine fachliche
Engineering-Freigabe ersetzen. Human Review kann fachliche Entscheidungen
autorisieren, aber keine technische Regression oder Syntaxvalidierung
ersetzen. Ebenso liefert eine externe Toolvalidierung keine Aussage über die
fachliche Korrektheit des zugrunde liegenden Modells.

Kapitel~7 stellt im nächsten Schritt die aus diesen V\&V-Aktivitäten
resultierenden Beobachtungen und Ergebnisse strukturiert dar. Auf dieser
Evidenzbasis werden die Ergebnisse anschließend in Kapitel~8 interpretiert
und zur Beantwortung der Forschungsfragen herangezogen.